\section{Introduction}

The order in which a compiler applies its optimization passes changes the code it produces.
Passes enable and disable each other. Inlining exposes constants that simplification folds, and
an early simplification can remove the pattern a later pass was designed to catch. Choosing a good
order is the phase-ordering problem, and the search space grows exponentially with sequence
length~\cite{almagor2004,kulkarni2012,ashouri2018survey}.

Recent work makes the search cheaper by summarising pairwise pass interactions in a graph.
DCO~\cite{dco2024} classifies ordered pass pairs by their measured dependence and combines short
sub-sequences along positive edges. Pan et al.~\cite{pan2025synergy} mine a synergy knowledge graph
of LLVM passes and use it to guide a genetic algorithm. Both build \emph{one} graph from a training
corpus and apply it to every program. This assumes that an interaction measured on other programs,
on unoptimized IR, still holds for a new program at the point in the pipeline where the passes
actually run.

We test that assumption directly in LLVM~\LLVMVersion. We measure \NPairs{} ordered pass-pair
interactions (\NPasses{} passes) on \NPrograms{} programs. For each program we do this at
\NSnapshots{} points along LLVM's own \Otwo{} pipeline, using \NExperiments{} compile-and-measure
experiments with no failures. We fixed four hypotheses, their tests and their thresholds in a
pre-registration before any full result existed. We report all four, including the one that
failed.

\paragraph{Contributions.}
\begin{itemize}
\item \textbf{Program-specific headroom (H1, supported).} With the same sequence builder and the same budget, ordering with
the program's own interaction graph instead of a static graph reduces IR size by \HOnePct\%. The reduction
carries over to machine code: \MSizeXPct\% smaller \texttt{.text} on x86-64 and \MSizeAPct\% on AArch64.
\item \textbf{Drift (H2, supported).} A program's interactions change as optimization proceeds. The
median rank correlation between interactions on raw IR and after the inliner and function
simplification is \HTwoRho. The largest change happens in the first simplification stage;
later stages resemble each other more.
\item \textbf{Predictability (H3, supported).} A per-program model trained on other programs removes
\HThreeGap\% of the static graph's error (residual $R^2$ = \HThreeRsq) and is better on
\HThreeBetter{} of \HThreePrograms{} held-out programs.
\item \textbf{A negative result (H4, not supported).} Re-predicting interactions along the
optimization trajectory does not beat the strongest search baseline at a budget of ten compilations.
Exploratory analysis locates the cause in the graph-based sequence builder, not in the predictions.
With a single compilation, however, predicted graphs do beat every search method.
\end{itemize}
The artifact will be made available.
